\section{Fixtures}

The positive RAIRE fixture is
\texttt{synthetic\_raire\_boundary\_package}. It is built by
\texttt{synthetic\_ranked\_choice\_boundary\_package} with method id
\texttt{raire-irv-v1}, run id \texttt{audit-run:raire-irv-boundary}, and
assorter id \texttt{raire-neb-not-eliminated-before-v1}. The synthetic reported
elimination order is \texttt{cand-c}, \texttt{cand-b}, \texttt{cand-a}. It has
one ranked-choice assertion, two ranked sample steps, and decision
\texttt{boundary}. The corresponding unit test verifies that the package passes
and preserves the elimination order.

The positive AWAIRE fixture is
\texttt{synthetic\_awaire\_boundary\_package}. It reuses the same helper and
ranked evidence but names method id \texttt{awaire-irv-v1}, run id
\texttt{audit-run:awaire-irv-boundary}, and assorter id
\texttt{awaire-adaptive-irv-v1}. Its test checks that the package verifies and
that the method id is the AWAIRE id.

The negative ranked-choice fixture is
\texttt{synthetic\_bad\_raire\_boundary\_package}. It starts from the RAIRE
boundary package and appends \texttt{cand-a} a second time to the first sampled
ranking. The ranked-choice verifier rejects the package as
\texttt{InvalidRankedChoiceSample}, proving that duplicate ranked choices are
not accepted as replayable boundary evidence.

These fixtures are intentionally small. They test method identity, elimination
order preservation, ranked sample shape, assertion kind, and malformed ranked
ballot rejection. They do not test IRV tabulation, RAIRE assertion generation,
assertion-level risk replay, or a real ranked CVR adapter.
